\section{Related Work}\label{sec:related}

\paragraph{Fixed sparse patterns.} Sparse Transformer
\citep{child2019sparse} introduced strided and fixed factorised patterns.
Longformer \citep{beltagy2020longformer} combined a sliding window with
task-specific global tokens, and BigBird \citep{zaheer2020bigbird} added
random attention and gave the resulting pattern a theoretical justification.
All three fix the pattern before training. Our local window and global summary
tokens are drawn directly from this line of work; the routed component is what
we add to it.

\paragraph{Kernel and low-rank approximations.} Linformer
\citep{wang2020linformer} projects keys and values to a fixed rank. Linear
attention \citep{katharopoulos2020linear} and Performer
\citep{choromanski2021performer} replace the softmax with a feature map,
recovering linear cost and a recurrent formulation. These methods approximate
the attention matrix; Corridor Attention computes an exact softmax over a
restricted support, which we prefer because the error is then structural and
inspectable rather than distributional.

\paragraph{Learned and content-based routing.} Routing Transformer
\citep{roy2021efficient} clusters queries and keys with online $k$-means and
attends within clusters. Reformer \citep{kitaev2020reformer} uses
locality-sensitive hashing to the same end. Both make the \emph{grouping}
content-dependent but leave each query's bandwidth as a consequence of cluster
occupancy, which is why both need balancing machinery. Corridor Attention fixes
the bandwidth per query and lets the routing vary, which we argue is the more
useful of the two invariants to hold fixed.

\paragraph{Retrieval-augmented and memory models.} A separate line of work
keeps long context outside the attention layer entirely, in a retrieval index
or a compressed memory \citep{rae2020compressive,borgeaud2022improving}. These
are complementary to our method rather than alternatives to it: Corridor
Attention is what we would use \emph{inside} such a model to read a long
retrieved passage.

